\ReviewerComment{The current introduction mainly emphasizes that long-term value function learning outperforms myopic planning methods, but it fails to elaborate on the practical dilemmas in existing satellite-ground scheduling research and engineering practice. Specifically, it does not explain why traditional scheduling methods such as MPC and dynamic programming are difficult to deploy in practical scenarios, nor does it distinguish between hard constraints and soft constraints in real satellite-ground scheduling systems. All experiments in this paper are conducted based on a soft-constraint clipping resource simulator, where the over-limit of resources such as energy and heat is only penalized without causing hardware failures. Therefore, it is necessary to supplement relevant explanations in the introduction and limitation sections: real satellite systems are dominated by hard constraints, where resource overrun will directly lead to task failure or hardware protection shutdown. The conclusions of this paper are only applicable to the soft-constraint simulation environment and cannot be directly transferred to real on-orbit satellite systems. It is recommended to add a dedicated paragraph to clarify the core differences between soft constraints in simulation and hard constraints in real satellite systems.}
\AuthorResponse{We added a dedicated Introduction paragraph and expanded the scope and limitations. The revised manuscript explains that clipping and penalties allow an episode to continue after a resource overrun, whereas operational systems can mask actions, reject tasks, or enter protective shutdown. We therefore delimit every principal claim to the synthetic soft-constraint benchmark.}
\RevisedExcerpt{A hard constraint changes the feasible action set and can make a policy with a lower penalized cost infeasible. Consequently, the ordering observed below is a statement about the specified soft-constraint benchmark and cannot be transferred directly to an on-orbit controller.}
\ManuscriptLocation{pp.~2--3, Introduction; p.~23, Section~5.1; p.~24, Conclusions.}

\ReviewerComment{The baseline algorithms adopted in this paper only include MPC-4, contextual bandit, simple threshold heuristic, and instantaneous cost minimization (argmin). Several representative mainstream algorithms in satellite edge scheduling are missing, including classic deep reinforcement learning (DRL) scheduling algorithms such as PPO and SAC, as well as traditional operation research baselines such as mixed integer programming (MIP) and rolling horizon heuristic scheduling. The current comparison only proves that the DQN family outperforms MPC-4, which is insufficient to demonstrate that DQN-based algorithms are the optimal algorithm set on the proposed benchmark. It is recommended to supplement the experimental results of at least 1--2 commonly used scheduling baselines (e.g., PPO) on this benchmark. If supplementary experiments cannot be completed due to computing resource limitations, the limitations of the current comparison should be clearly stated in the discussion section, and the potential performance of other alternative algorithms should be further discussed.}
\AuthorResponse{We added an eight-step greedy rolling-look-ahead baseline under the same perfect forecast used by MPC-4. We also added verified literature discussion of PPO, SAC, mixed-integer programming, and dynamic programming, and narrowed the claim to the implemented comparator set. The revised manuscript states that these alternatives may outperform DQN under different hard-feasibility, control-space, or global-look-ahead conditions. The new baseline is reported on pp.~20--21 (Table~12).}
\RevisedExcerpt{The greedy rollout reduced mean cost relative to MPC-4 (88.79 versus 90.83) but remained above all six DQN means (77.66--78.40). Every paired DQN--greedy difference was negative, ranging from $-11.12$ to $-10.39$; all 95\% bootstrap intervals excluded zero and the largest Holm-adjusted sign-test value was 0.000011.}
\ManuscriptLocation{pp.~2--3, Related Work; pp.~20--21, Section~4.6 and Table~12; p.~23, Section~5.2.}

\ReviewerComment{Section 3.4.1 clarifies that the auxiliary loss weight $\lambda=0.3$ is only tuned for the centered full-action model, and all other DQN variants directly adopt this hyperparameter without independent optimization. Such experimental setting makes it impossible to decouple the individual contributions of the centralization mechanism, Bellman bootstrap strategy, and auxiliary loss coefficient to the overall performance. The results in Table 7 only reflect the combined effect of multiple factors and cannot explain the internal optimization mechanism of the algorithm. Supplementary ablation experiments are suggested to explore the independent effectiveness of each module and hyperparameter.}
\AuthorResponse{We agree with the identifiability concern. A separate grid search for every auxiliary objective is outside the present revision scope. We therefore recast these rows as fixed-coefficient descriptive comparisons, explain the scale mismatch among auxiliary losses, and remove causal component claims. The manuscript now states that factorial identification requires objective-specific tuning and controlled gradient normalization.}
\RevisedExcerpt{The uncentered full-action, immediate-advantage, and Double DQN auxiliary variants inherited $\lambda=0.3$ without method-specific calibration. Because the auxiliary losses differ in scale and error structure, these are fixed-coefficient implementation comparisons.}
\ManuscriptLocation{p.~11, Section~3.4.1; pp.~17--18, Tables~7--8; p.~23, Sections~5.1--5.2.}

\ReviewerComment{The MPC-$H$ planner adopted in this paper possesses perfect forecast information of all future tasks and link states in the subsequent $H$ steps, while the DQN policy can only observe the current system state without access to future task information. The inconsistent information acquisition conditions between the two algorithms artificially raise the performance upper bound of the MPC baseline. In practical engineering scenarios, accurate full prediction of future tasks is hardly achievable. It is necessary to highlight the information asymmetry between the compared algorithms in the experimental illustration. Additionally, MPC with imperfect prediction (i.e., prediction noise) can be added as a supplementary baseline. If no new experiments are supplemented, the discussion section should analyze the performance degradation of MPC and the corresponding changes in the advantages of DQN family algorithms when future task prediction errors exist.}
\AuthorResponse{We now state this asymmetry in Methods, Results captions, and Discussion. The new greedy rollout deliberately receives the same forecast as MPC. We did not add a noisy-MPC experiment because a single arbitrary noise distribution would not represent operational forecast quality. Instead, we discuss calibrated prediction distributions and structural link-state errors as necessary future tests. The present comparison is conservative with respect to DQN because only the planners receive future tasks and link states.}
\RevisedExcerpt{Both MPC and greedy rollout receive perfect next-task and link-state previews; learned policies receive only the current state, so the comparison includes an explicit information asymmetry.}
\ManuscriptLocation{pp.~12--13, Section~3.4.4; p.~16, Figure~5 caption; p.~21, Section~4.6; p.~22, Discussion.}
